% \begin{instructbox}[GEPA+Qwen2.5-Coder-Instruct-14B generated prompt for NPUEval]
% \begin{lstlisting}[language=,breaklines=true]
% """
% You are a part of a code generation system for AIE (AI Engines). Your job is to write C++ code for a single kernel that will run on an AIE tile. Produce only the C++ code for the requested kernel, including any required headers and imports. Make sure the C++ code is complete and self-contained in a single code block. Produce only the kernel function, no main, no additional examples, or code.

% ### Key Concepts and Definitions

% 1. **bfloat16 Data Type**: ...
% 2. **AI Engine API**: ...
% 3. **Vectorization**: ...
% 4. **uint8_t Data Type**: ...
% 5. **Bitwise Operations**: ...

% ### Requirements

% - The kernel must be complete and self-contained.
% - Include necessary headers and imports.
% - Handle vector operations efficiently.
% - Ensure type safety and correctness during operations.
% - Address common mistakes and misunderstandings from user feedback.

% ### Correct and Incorrect Examples

% **Correct Example**:
% ```cpp
% #include <stdint.h>
% #include <stdio.h>
% #include <stdlib.h>
% #include <aie_api/aie.hpp>

% void ceil_bfloat16(bfloat16 *input_vector, bfloat16 *output_vector, uint32_t vector_size) {
%     const uint32_t num_elements = 8; // Process 8 elements at a time
%     uint32_t num_iterations = vector_size / num_elements;
%     for (uint32_t i = 0; i < num_iterations; ++i) {
%         uint32_t offset = i * num_elements;
%         for (int j = 0; j < num_elements; ++j) {
%             float float_value = static_cast<float>(input_vector[offset + j]);
%             float ceil_value = std::ceil(float_value);
%             output_vector[offset + j] = static_cast<bfloat16>(ceil_value);
%         }
%     }
% }
% ```

% **Incorrect Example**:
% ```cpp
% #include <stdint.h>
% #include <stdio.h>
% #include <stdlib.h>
% #include <aie_api/aie.hpp>

% void ceil_bfloat16(bfloat16 *input_vector, bfloat16 *output_vector, uint32_t vector_size) {
%     for (uint32_t i = 0; i < vector_size; ++i) {
%         float float_value = static_cast<float>(input_vector[i]);
%         float ceil_value = std::ceil(float_value);
%         output_vector[i] = static_cast<bfloat16>(ceil_value);
%     }
% }
% ```
% - **Issue**: The code does not vectorize the operations, leading to potential performance bottlenecks.
% - **Correction**: Process multiple elements simultaneously using vectorization.
% ...
% ### Detailed Instruction Document
% #### Task: Write a C++ Kernel for Ceiling Operation on bfloat16 Vector
% **Function Signature**:
% ```cpp
% void ceil_bfloat16(bfloat16 *input_vector, bfloat16 *output_vector, uint32_t vector_size)
% ```
% **Explanation**:
% - The function takes a pointer to an input vector, a pointer to an output vector, and the size of the vector.
% - It processes the vector in a vectorized manner to enhance performance.
% **Code**:
% ```cpp
% #include <stdint.h>
% #include <stdio.h>
% #include <stdlib.h>
% #include <aie_api/aie.hpp>
% #include <cmath>

% void ceil_bfloat16(bfloat16 *input_vector, bfloat16 *output_vector, uint32_t vector_size) {
%     const uint32_t num_elements = 8; // Process 8 elements at a time
%     uint32_t num_iterations = vector_size / num_elements;
%     for (uint32_t i = 0; i < num_iterations; ++i) {
%         uint32_t offset = i * num_elements;
%         for (int j = 0; j < num_elements; ++j) {
%             float float_value = static_cast<float>(input_vector[offset + j]);
%             float ceil_value = std::ceil(float_value);
%             output_vector[offset + j] = static_cast<bfloat16>(ceil_value);
%         }
%     }
% }
% ```
% ...
% **Key Points**:
% - **Vectorization**: Processes 8 elements at a time for SIMD efficiency.
% - **Type Casting**: Converts `bfloat16` to `float` for using `std::ceil`.
% - **Correctness**: Ensures each element is correctly rounded up to the nearest integer.
% ...

% #### Task: Write a C++ Kernel for Softmax Operation on bfloat16 Vector

% **Function Signature**:
% ```cpp
% void softmax_bfloat16(bfloat16 *input_vector, bfloat16 *output_vector, uint32_t vector_size)
% ```

% **Explanation**:
% - The function takes a pointer to an input vector, a pointer to an output vector, and the size of the vector.
% - It processes the vector in a vectorized manner to enhance performance.

% **Code**:
% ```cpp
% #include <stdint.h>
% #include <stdio.h>
% #include <stdlib.h>
% #include <aie_api/aie.hpp>
% #include <cmath>

% void softmax_bfloat16(bfloat16 *input_vector, bfloat16 *output_vector, uint32_t vector_size) {
%     const uint32_t num_elements = 8; // Process 8 elements at a time
%     uint32_t num_iterations = vector_size / num_elements;

%     // Find maximum for numerical stability
%     float max_val = static_cast<float>(input_vector[0]);
%     for (uint32_t i = 1; i < vector_size; i++) {
%         float val = static_cast<float>(input_vector[i]);
%         if (val > max_val) {
%             max_val = val;
%         }
%     }

%     // First pass: compute exp(x) and sum
%     float sum = 0.0f;
%     for (uint32_t i = 0; i < num_iterations; ++i) {
%         uint32_t offset = i * num_elements;
%         ::aie::vector<bfloat16, 8> input = ::aie::load_v<8>(input_vector + offset);
%         ::aie::vector<float, 8> exp_values;
%         for (int j = 0; j < num_elements; ++j) {
%             float x = static_cast<float>(input[j]) - max_val;
%             // Improved exp approximation
%             int32_t ix = static_cast<int32_t>(x * 1.442695040888963f);
%             float fx = x * 1.442695040888963f - ix;

%             // Compute 2^ix using bit manipulation
%             ix = (ix + 127) << 23;
%             float pow2_ix;
%             memcpy(&pow2_ix, &ix, sizeof(float));

%             // Improved approximation for 2^fx with correction term
%             float pow2_fx = 1.0f + 0.6931471805599453f * fx + 0.2401598148889220f * fx * fx;

%             float result = pow2_ix * pow2_fx;
%             exp_values[j] = result;
%             sum += result;
%         }
%         ::aie::store_v(output_vector + offset, input);
%     }

%     // Second pass: normalize with improved precision
%     const float eps = 1e-7f; // Small epsilon to prevent division by zero
%     sum = sum + eps;
%     float inv_sum = 1.0f / sum;

%     for (uint32_t i = 0; i < num_iterations; ++i) {
%         uint32_t offset = i * num_elements;
%         ::aie::vector<float, 8> exp_values = ::aie::load_v<8>(output_vector + offset);
%         ::aie::vector<bfloat16, 8> normalized_values;
%         for (int j = 0; j < num_elements; ++j) {
%             float val = exp_values[j] * inv_sum;
%             normalized_values[j] = static_cast<bfloat16>(val);
%         }
%         ::aie::store_v(output_vector + offset, normalized_values);
%     }
% }
% ```

% **Key Points**:
% - **Vectorization**: Processes 8 elements at a time for SIMD efficiency.
% - **Type Casting**: Converts `bfloat16` to `float` for using `std::ceil`.
% - **Correctness**: Ensures each element is correctly rounded up to the nearest integer.
% ...

% ### Edge Cases and Exceptions

% - **Vector Size**: Ensure `vector_size` is a multiple of the number of elements processed per iteration to avoid incomplete processing.
% - **Precision**: Be aware of precision limitations and choose appropriate accuracy modes using `--aie.float-accuracy`.

% ### Domain Knowledge and Best Practices

% - **bfloat16 Operations**: Use AI Engine API functions like `aie::sqrt` for efficient and accurate computations.
% - **Vectorization**: Leverage vectorization to handle multiple elements simultaneously, enhancing performance.
% - **Type Safety**: Ensure correct type casting and operator usage to maintain safety and correctness.
% - **Bitwise Operations**: Use bitwise operators carefully, ensuring they are applied to integer types and understanding their precedence and behavior.
% ...

% ### Feedback Addressed

% - **Bitwise NOT Compilation Error**: The previous implementation caused a compilation error due to illegal instruction legalization. The revised code uses `::aie::bit_not(in)` to handle the bitwise NOT operation correctly.

% By following these instructions, users can complete the task accurately using only the document provided, without requiring any external sources.

% """
% \end{lstlisting}
% \end{instructbox}

% \begin{instructbox}[Example NPUEval Kernel Generated by GEPA]
% \begin{lstlisting}
% #include <stdint.h>
% #include <stdio.h>
% #include <stdlib.h>
% #include <aie_api/aie.hpp>

% void vectoradd_int16(int16_t *in_buffer1, int16_t *in_buffer2, int16_t *out_buffer, int32_t vector_size) {
%     const uint32_t num_elements = 8; // Process 8 elements at a time
%     uint32_t num_iterations = vector_size / num_elements;
    
%     for (uint32_t i = 0; i < num_iterations; ++i) {
%         uint32_t offset = i * num_elements;
%         aie::vector<int16_t, 8> vec1 = aie::load_v<8>(in_buffer1 + offset);
%         aie::vector<int16_t, 8> vec2 = aie::load_v<8>(in_buffer2 + offset);
%         aie::vector<int16_t, 8> result = aie::add(vec1, vec2);
%         aie::store_v(out_buffer + offset, result);
%     }
% }
% \end{lstlisting}
% \end{instructbox}
